\section{Conclusion and Future Work}
\label{sec:conclusion}

In this work, we have mathematically deconstructed and empirically resolved the core limitations of existing test-time ensembling and model merging paradigms. By exposing that standard parametric and wave-inspired routers are fundamentally susceptible to few-shot overfitting (the Dynamic Routing Paradox) and sequential stream instability (Vectorization Collapse), we highlighted the urgent need for mathematically robust ensembling alternatives.

To resolve these challenges, we introduced \textbf{Fisher-Information Optimal Subspace Routing (FIOSR)}, a completely training-free and parameter-free dynamic ensembling framework that treats the parameter space as a Riemannian manifold. By modeling local parameter sensitivities using a smoothed and power-scaled diagonal empirical Fisher Information Matrix, we formulated an analytical local Riemannian metric tensor. This tensor warps the coordinate space during subspace projection, suppressing noise and task-irrelevant coordinates while highlighting highly discriminative dimensions. Integrated with Class-Size Scaling Calibration (CSC) and Micro-Batch Homogenization (MBH), FIOSR completely bypasses test-time parameter optimization. It achieves near-perfect specialized routing, recovering the oracle performance ceilings of individual experts ($\approx 100\%$ on MNIST and FashionMNIST, and near-ceiling accuracy on CIFAR-10 and SVHN) while maintaining absolute ensembling stability across all stream batching regimes ($B=1$ to $512$).

\subsection{System-Level Trade-offs and Gating Safeguards}
While MBH elegantly prevents \emph{heterogeneity collapse}, ensembling mixed streams in real-world deployments introduces crucial systems-level trade-offs. Partitioning a heterogeneous batch of size $B$ into $G \le K$ micro-batches requires running $G$ separate sequential forward passes on $G$ dynamically merged models, adding non-trivial execution latency and memory bandwidth overhead. While our evaluation demonstrates that Top-$M$ Expert Gating with $M=1$ (hard Top-1 routing) completely eliminates sequential MBH overhead (restricting active forward passes to 1) while achieving an exceptional joint accuracy of \textbf{76.87\%} ($+8.84\%$ over flat Cosine), this setup represents a conceptual compromise: it collapses ensembling back to a hard task-routing selection mechanism, losing the benefits of weight-space ensembling for that sample in exchange for absolute computational efficiency. In production pipelines, practitioners must choose where to stand on this latency-accuracy curve, balancing the rich multi-task ensembling benefits of larger $M$ against the absolute single-pass execution speed of $M=1$.

\subsection{Roadmap for End-to-End Physical Evaluation}
To transition FIOSR from simulated coordinate spaces to end-to-end physical deployments on actual large language models (e.g., LLaMA-3-8B) or Vision Transformers (e.g., ViT-Base), we outline a concrete four-step evaluation roadmap:
\begin{enumerate}\setlength{\itemsep}{0.5pt}\setlength{\parskip}{0pt}
    \item \textbf{Activation Collection \& Support Splits:} First, fine-tune $K$ specialized expert adapters (e.g., using LoRA) on target domains (such as distinct GLUE tasks or specialized image classification tasks). For each task, pass a microscopic calibration split of $N_c=16$ samples through the shared pre-trained base model backbone, extracting intermediate hidden representations (e.g., outputs of the multi-layer perceptron or self-attention layers).
    \item \textbf{On-The-Fly FIM Estimation:} Second, calculate the pooled within-class covariance and diagonal representation-space Fisher Information Matrix (dFIM) directly on these activation vectors. Since these activations are low-dimensional (typically $d \in [768, 4096]$) compared to full parameter counts, computing class-conditional coordinate variances takes only a few milliseconds, making on-the-fly metric warp construction computationally negligible.
    \item \textbf{Online Metric Warping and Routing:} Third, during test-time inference on non-stationary, heterogeneous streams, use our smoothed dFIM metric tensor to warp activation coordinates and project test samples onto class prototype vectors using Fisher-Weighted Cosine Similarity. Apply CC-CSC to calibrate vocabularies and MBH to partition batches.
    \item \textbf{Dynamic Adapter Merging:} Fourth, utilize libraries like Hugging Face PEFT or custom CUDA kernels (e.g., Punica) to dynamically merge the LoRA weights of active experts on the fly based on the computed routing coefficients, performing single-pass ensembling for homogeneous sub-batches.
\end{enumerate}

\subsection{Future Directions} From a theoretical perspective, several compelling paths lie ahead:
\begin{itemize}\setlength{\itemsep}{0.5pt}\setlength{\parskip}{0pt}
    \item \textbf{Non-Diagonal Riemannian Metrics:} While the diagonal FIM offers high computational efficiency, future research should explore block-diagonal or low-rank approximations of the Fisher Information Matrix (such as K-FAC \cite{martens2015optimizing}) to capture complex, cross-dimensional parameter correlations.
    \item \textbf{Active Online Riemannian Adaptation:} Under non-stationary environments with continuous domain drift, incorporating control-theoretic Lyapunov stability constraints into the Fisher estimation process will enable real-time, provably stable metric tensor updates.
\end{itemize}
Ultimately, we hope this work inspires researchers to move away from flat Euclidean ensembling heuristics and embrace the rich, rigorous foundations of information-geometric and Riemannian manifold analysis in modular deep learning.
